\section*{अध्याय \ref{ch:g4:fractions} --- पहली भिन्नें}

\begin{solution}{exo:g4:fractions:1}
पिज़्ज़ा: $\frac56$ बचा। चॉकलेट: $\frac38$ खाई गई। झंडा:
$\frac13$ रंगा।
\end{solution}

\begin{solution}{exo:g4:fractions:2}
$\frac35$ तीन हिस्से रँगता है, $\frac25$ केवल दो: $\frac35$ बड़ी
रँगाई है।
\end{solution}

\begin{solution}{exo:g4:fractions:3}
$18$ का \omterm{def:g3:division:half}{आधा}: $9$। $20$ की \omterm{def:g3:division:half}{चौथाई}: $5$। $21$ की तिहाई: $7$।
$20$ का $\frac34$: तीन \omterm{def:g3:division:half}{चौथाई}, $3 \times 5 = 15$।
\end{solution}

\begin{solution}{exo:g4:fractions:4}
$\frac13$: एक कदम; $\frac33$: $1$ पर; $\frac53$: $1$ से दो कदम
आगे। ठीक $2$ पर आने वाली \omterm{def:g4:fractions:def}{भिन्न} $\frac63$ है।
\end{solution}

\begin{solution}{exo:g4:fractions:5}
$\frac38 < \frac68$; \quad $\frac55 = 1$; \quad $\frac74 > 1$;
\quad $\frac12 = \frac24$।
\end{solution}

\begin{solution}{exo:g4:fractions:6}
$\frac26 + \frac36 = \frac56$; \quad
$\frac58 - \frac28 = \frac38$; \quad
$\frac14 + \frac24 = \frac34$; \quad
$1 = \frac33$, इसलिए $1 - \frac13 = \frac23$।
\end{solution}

\begin{solution}{exo:g4:fractions:7}
$\frac54 = 1 + \frac14$; \quad
$\frac73 = 2 + \frac13$; \quad
$\frac92 = 4 + \frac12$।
\end{solution}

\begin{solution}{exo:g4:fractions:8}
आधा घंटा: $30$ मिनट। \omterm{def:g3:division:half}{चौथाई}: $15$ मिनट। तीन \omterm{def:g3:division:half}{चौथाई}: $45$
मिनट। तिहाई: $20$ मिनट।
\end{solution}

\begin{solution}{exo:g4:fractions:9}
$24$ का $\frac14$ यानी $6$ विद्यार्थी चश्मा पहनते हैं;
$24 - 6 = 18$ नहीं पहनते।
\end{solution}

\begin{solution}{exo:g4:fractions:10}
मिलकर: $\frac38 + \frac48 = \frac78$। बची: $\frac18$। सैम ने
ज़्यादा खाया ($4$ आठवें भाग बनाम $3$)।
\end{solution}

\begin{solution}{exo:g4:fractions:11}
हाँ, ऐसा हो सकता है: छोटे पिज़्ज़े का \omterm{def:g3:division:half}{आधा}, बहुत बड़े पिज़्ज़े की
\omterm{def:g3:division:half}{चौथाई} से कम खाना हो सकता है। भिन्नें \emph{एक ही पूरे} के
हिस्सों की तुलना करती हैं; ``किसी चीज़'' के $\frac12$ और
$\frac14$ की तुलना से पहले यह पता होना चाहिए कि दोनों चीज़ें बराबर हैं।
\end{solution}
